
\documentclass{article}
\usepackage[utf8]{inputenc}
\usepackage{amsmath}
\usepackage{amsfonts}
\usepackage{amssymb}
\usepackage{amsthm}
\usepackage{boxedminipage}
\usepackage{booktabs}
\usepackage{tabularx}
\newtheorem{theorem}{Theorem}[section]
\emergencystretch=2em

\title{Master Manuscript V4: Structural Consistency and Residue Sanity}
\author{Thomas Hoffbauer}
\date{August 16, 2026}

\begin{document}

\maketitle

\section{The Residue Sanity Gate: Metric-Volumetric Invariance}

To ensure structural consistency in the EABC framework, we establish the Residue Sanity Gate. For any full-dimensional Euclidean lattice $\Lambda \subset \mathbb{R}^d$, the radial Epstein-Zeta function is defined as:
\[
Z_{\Lambda}^{\mathrm{len}}(s) = \sum_{0 \neq x \in \Lambda} \|x\|^{-s}.
\]
The fundamental geometric consistency condition at the pole $s=d$ is:
\begin{equation}
\boxed{\operatorname*{Res}_{s=d} Z_{\Lambda}^{\mathrm{len}}(s) = \frac{2\pi^{d/2}}{\Gamma(d/2)\operatorname{covol}(\Lambda)}}
\end{equation}

\noindent\textbf{Crucial Clarification on Invariance:} 
The residue is \emph{not} a topological invariant, but a \emph{metric-volumetric lattice invariant}. Under dilation $\Lambda \mapsto c\Lambda$, the zeta function transforms as $Z_{c\Lambda}^{\mathrm{len}}(s) = c^{-s}Z_\Lambda^{\mathrm{len}}(s)$, yielding:
\[
\operatorname*{Res}_{s=d} Z_{c\Lambda}^{\mathrm{len}} = c^{-d} \operatorname*{Res}_{s=d} Z_\Lambda^{\mathrm{len}}.
\]
More generally, for any $A \in GL(d, \mathbb{R})$, we have $\operatorname{covol}(A\Lambda) = |\det A|\operatorname{covol}(\Lambda)$, leading to:
\[
\boxed{
\operatorname*{Res}_{s=d} Z_{A\Lambda}^{\mathrm{len}} = |\det A|^{-1} \operatorname*{Res}_{s=d} Z_\Lambda^{\mathrm{len}}.
}
\]
Thus, the residue is invariant under all volume-preserving linear transformations ($|\det A| = 1$), including orientation-reversing ones.

\section{Divisor Information and Winding Audit (Argument Principle)}

For a meromorphic function, the winding number along a closed contour $\gamma$ in the resolvent set is given by the argument principle:
\[
\nu_\gamma = \frac{1}{2\pi i} \oint_\gamma \frac{Z'(s)}{Z(s)}\,ds = N_Z(\gamma) - N_P(\gamma) \in \mathbb{Z}.
\]
This integer is invariant under homotopies of the contour within the complement of the divisor. This forms the clean dichotomy:
\[
\begin{array}{rcl}
\operatorname*{Res}_{s=d} Z_\Lambda &\longleftrightarrow& \text{metric/volumetric pole strength}, \\[2mm]
\displaystyle \frac{1}{2\pi i}\oint_\gamma d\log Z_\Lambda &\longleftrightarrow& \text{discrete divisor / winding information}.
\end{array}
\]

\section{Comparative Divisor and Winding Audit for $D_4$ and $E_8$}

Because $D_4$ lives in dimension 4 and $E_8$ in dimension 8, there is no canonical continuous cross-dimensional transition along which a single winding number could be transported directly from $D_4$ to $E_8$. Instead, we perform a rigorous comparative divisor audit by analyzing the divisors $\operatorname{div} Z_{D_4}^{\mathrm{len}}$ and $\operatorname{div} Z_{E_8}^{\mathrm{len}}$ using concentric contours centered at the respective dimensional poles:
\[
\gamma_{D_4}(R): |s-4|=R, \quad \gamma_{E_8}(R): |s-8|=R.
\]
The audit maps the discrete step functions:
\[
R \mapsto \bigl(\nu_{D_4}(R), \nu_{E_8}(R)\bigr),
\]
designed to determine the radial distances at which zeros and secondary poles appear and to identify cancellations among factor singularities. No cross-dimensional transport law is assumed.



\section{Monomial Fourier--Projection Uniqueness Theorem}

Define the $L=1$ projection functional
\[
\mathcal P[f] := 2\int_0^1 (1-s)f(s)\,ds
\]
and, for integers $a\ge 0$ and $m\ge 1$, the monomial Fourier family
\[
f_{a,m}(s):=s^a\sin(2\pi m s).
\]
Set $k=2\pi m$ and
\[
M_a(k):=\int_0^1 s^a\sin(ks)\,ds.
\]
For $a\ge2$, two integrations by parts, together with $\sin k=0$ and $\cos k=1$, give
\begin{equation}
M_a(k)=-\frac1k-\frac{a(a-1)}{k^2}M_{a-2}(k),
\end{equation}
with initial values $M_0(k)=0$ and $M_1(k)=-1/k$.

Introduce
\[
x:=\frac1{k^2},\qquad J_a(x):=kM_a(k).
\]
Then
\[
J_0=0,\qquad J_1=-1,
\]
and
\begin{equation}
J_a(x)=-1-a(a-1)xJ_{a-2}(x),\qquad a\ge2.
\end{equation}
Hence $J_a\in\mathbb Z[x]\subset\mathbb Q[x]$ for every $a$.

Since
\[
\mathcal P[f_{a,m}]=2(M_a-M_{a+1}),
\]
define
\[
D_a(x):=J_a(x)-J_{a+1}(x).
\]
Then
\begin{equation}
\mathcal P[f_{a,m}]=\frac{2}{k}D_a\!\left(\frac{1}{4\pi^2m^2}\right).
\end{equation}

\begin{theorem}[Monomial Fourier--Projection Uniqueness]
For every $a\in\mathbb N_0$ and $m\in\mathbb N$,
\[
\boxed{\mathcal P\!\left[s^a\sin(2\pi m s)\right]=0\iff a=1.}
\]
\end{theorem}

\begin{proof}
For $a=1$, the recursion gives $J_1=J_2=-1$, hence $D_1\equiv0$.
For $a=0$, $D_0=1$. For $a=2$, one finds $D_2(x)=-6x\not\equiv0$.
For every $a\ge3$, the recursion yields
\[
J_a(x)=-1+a(a-1)x+O(x^2),
\qquad
J_{a+1}(x)=-1+a(a+1)x+O(x^2),
\]
so
\[
D_a(x)=-2ax+O(x^2),
\]
and therefore $D_a$ is a nonzero polynomial in $\mathbb Q[x]$.

For fixed $m\ge1$,
\[
x_m=\frac1{4\pi^2m^2}
\]
is transcendental because $\pi$ is transcendental. A nonzero polynomial with rational coefficients has only algebraic roots. Consequently $D_a(x_m)=0$ can occur only when $D_a\equiv0$, which happens exactly for $a=1$.
\end{proof}

The cancellation at $a=1$ is nontrivial at the level of moments:
\[
M_1(k)=M_2(k)=-\frac1k=-\frac1{2\pi m}\neq0,
\]
while
\[
\mathcal P[s\sin(2\pi m s)]=2(M_1-M_2)=0.
\]
Thus the linear exponent is the unique kernel exponent within this monomial Fourier family. This statement is algebraic-analytic; no physical vacuum interpretation is implied by the theorem itself.



% -----------------------------------------------------------------------------
% Intended placement in the extended master architecture: Section 4.5.2.
% The present compact V4 source does not yet contain the complete 4.5 hierarchy;
% therefore the semantic number is retained in the heading without falsifying the
% automatic section counters of this standalone source.
\subsection*{4.5.2 Disjoint Logarithmic Shell Audit: Support Weights and Detailed Balance}
\addcontentsline{toc}{subsection}{4.5.2 Disjoint Logarithmic Shell Audit: Support Weights and Detailed Balance}

To separate genuine bulk behavior from smoothing caused by cumulative averages, we
repeat the support-weight and reversibility audit on disjoint prime-value shells.  Let
\(\widehat\pi^{\mathrm{shell}}=(\widehat\pi_2,\widehat\pi_3,\widehat\pi_4)\)
denote the empirical frequencies of the three coarse-grained support sectors
\(S_2,S_3,S_4\) in a shell, and let
\[
\varpi=\left(\frac12,\frac14,\frac14\right)
\]
be the predicted reference weight vector.  For the shell transition matrix
\(T^{\mathrm{shell}}\), define the detailed-balance defect by
\begin{equation}
D_{\mathrm{DB}}^{\mathrm{shell}}
:=
\left\|
\operatorname{diag}(\varpi)T^{\mathrm{shell}}
-
\bigl(T^{\mathrm{shell}}\bigr)^{\!T}\operatorname{diag}(\varpi)
\right\|_F .
\end{equation}
Thus \(D_{\mathrm{DB}}^{\mathrm{shell}}=0\) is the exact detailed-balance condition
for the chosen reference weights, whereas a nonzero value measures the finite-shell
defect in Frobenius norm.

\begin{table}[htbp]
\centering
\small
\caption{Disjoint logarithmic shell audit for primes up to \(p=500{,}000\).  The
reported values are finite-range numerical data.  They provide evidence for
stabilization of the support frequencies and a reduced detailed-balance defect in the
largest shell, but do not by themselves prove an asymptotic limit.}
\label{tab:disjoint-shell-audit}
\begin{tabularx}{\textwidth}{@{}c c r c c X@{}}
\toprule
Shell & Prime-value range \(p\) & Count &
\(\widehat\pi^{\mathrm{shell}}(S_2,S_3,S_4)\) &
\(D_{\mathrm{DB}}^{\mathrm{shell}}\) &
Diagonal entries \((T_{22},T_{33},T_{44})\) \\
\midrule
1 & \([100,1{,}000)\) & 143 &
\((0.4825,0.2587,0.2587)\) & 0.035057 &
\((0.3529,0.0541,0.1622)\) \\
2 & \([1{,}000,5{,}000)\) & 501 &
\((0.4970,0.2475,0.2555)\) & 0.008821 &
\((0.3735,0.1371,0.1181)\) \\
3 & \([5{,}000,20{,}000)\) & 1,593 &
\((0.4997,0.2530,0.2473)\) & 0.034353 &
\((0.4138,0.1241,0.1269)\) \\
4 & \([20{,}000,100{,}000)\) & 7,330 &
\((0.4991,0.2508,0.2502)\) & 0.008211 &
\((0.4155,0.1447,0.1468)\) \\
5 & \([100{,}000,500{,}000)\) & 31,946 &
\((0.4992,0.2507,0.2501)\) & \(\mathbf{0.006970}\) &
\((0.4225,0.1587,0.1576)\) \\
\bottomrule
\end{tabularx}
\end{table}

\paragraph{Support-weight stabilization.}
The disjoint-shell frequencies approach the reference vector closely over the tested
range.  In the largest shell,
\[
\widehat\pi^{\mathrm{shell}}=(0.4992,0.2507,0.2501),
\]
so the componentwise deviations from \(\varpi\) are
\[
(-8\times10^{-4},\;7\times10^{-4},\;1\times10^{-4}).
\]
This is strong finite-range shell evidence that the observed support weights are stable
near the predicted values.  It is not, on its own, a proof that
\(\widehat\pi^{\mathrm{shell}}\to\varpi\) as \(p\to\infty\).  Such an asymptotic
statement requires an independent analytic argument, including an exact description of
the pushforward from the underlying residue classes to the three support sectors.

\paragraph{Detailed-balance defect.}
The sequence
\[
0.035057,\;0.008821,\;0.034353,\;0.008211,\;0.006970
\]
is manifestly non-monotone.  Nevertheless, the largest shell attains the smallest
reported defect.  The appropriate finite-data conclusion is therefore an overall
endpoint reduction accompanied by substantial shell-to-shell fluctuations.  This is
compatible with finite-size arithmetic correlations, but the present audit does not
identify Hardy--Littlewood correlations as the unique cause and does not establish
\[
\lim_{p\to\infty}D_{\mathrm{DB}}^{\mathrm{shell}}=0.
\]

\paragraph{Diagonal persistence.}
The diagonal entries are reported without qualitative dynamical labels.  In the largest
shell one finds
\[
T_{22}=0.4225,\qquad T_{33}=0.1587,\qquad T_{44}=0.1576.
\]
Their differing magnitudes are an empirical property of the implemented coarse-grained
transition matrix.  Any combinatorial explanation in terms of preimage multiplicities or
cross-transition counts must be established separately from the numerical table.

\paragraph{Epistemic status.}
The shell audit strengthens the numerical evidence because the intervals are disjoint and
therefore do not inherit the deterministic smoothing of cumulative averages.  The current
status is
\[
\boxed{
\begin{array}{ll}
\text{Support-weight convergence} & \mathrm{STRONG\ FINITE\ RANGE\ EVIDENCE},\\
\text{Detailed-balance defect} & \mathrm{ENDPOINT\ REDUCTION\ WITH\ FLUCTUATIONS},\\
\text{Asymptotic detailed balance} & \mathrm{NOT\ YET\ PROVED},\\
\text{Borel-measure generation} & \mathrm{NOT\ CLAIMED}.
\end{array}}
\]
The terminology ``Borel weights'' is therefore avoided here: \(\varpi\) is treated as a
predicted support-weight vector (or, once proved, as a Dirichlet-pushforward weight),
while the standard Borel measure-theoretic framework is presupposed rather than derived
from the discrete arithmetic model.


\section{Residue Specializations for $E_8$ and $D_4$}

For the standard $E_8$ lattice, $\operatorname{covol}(E_8)=1$, and therefore
\[
\boxed{\operatorname*{Res}_{s=8}Z_{E_8}^{\mathrm{len}}(s)=\frac{\pi^4}{3}.}
\]
For the standard root lattice
\[
D_4=\{x\in\mathbb Z^4: x_1+x_2+x_3+x_4\equiv0\pmod2\},
\]
we have $\operatorname{covol}(D_4)=2$, hence
\[
\boxed{\operatorname*{Res}_{s=4}Z_{D_4}^{\mathrm{len}}(s)=\pi^2.}
\]
If $D_4$ is rescaled to covolume $1$, the corresponding residue is $2\pi^2$.
These values are normalization-sensitive consequences of the general residue formula and constitute mandatory consistency checks for any proposed product or Dirichlet-series representation.

\subsection{Status of the comparative winding stairs}

The comparative winding functions are defined by
\[
\nu_{D_4}(R)=\frac{1}{2\pi i}\oint_{|s-4|=R}d\log Z_{D_4}^{\mathrm{len}}(s),
\qquad
\nu_{E_8}(R)=\frac{1}{2\pi i}\oint_{|s-8|=R}d\log Z_{E_8}^{\mathrm{len}}(s).
\]
Their step locations are determined by the complete divisors of the corresponding meromorphic continuations. A proposed factorization may be used for this purpose only after it has passed both the pole-structure check and the Residue Sanity Gate. Accordingly, the comparative staircase is defined here as an audit object; a full $D_4$ staircase is not claimed until an independently validated $D_4$ factorization or divisor description is supplied.

\section{Claim-Wall}

\begin{itemize}
    \item \textbf{Metric-Volumetric Residue Formula:} PROVEN
    \item \textbf{$E_8$ specialization ($\operatorname{covol}=1$) and standard $D_4$ specialization ($\operatorname{covol}=2$):} PROVEN
    \item \textbf{Argument Principle / Winding Invariance:} PROVEN
    \item \textbf{Konkrete Produktdarstellung:} VALIDATED (only if pole structure matches the Gate)
    \item \textbf{Comparative Divisor Audit ($D_4$ vs $E_8$):} DEFINED / AUDIT PENDING
    \item \textbf{Monomial Fourier--Projection Uniqueness Theorem:} PROVEN
    \item \textbf{Disjoint-shell support-weight stabilization:} STRONG FINITE-RANGE EVIDENCE
    \item \textbf{Detailed-balance shell defect:} ENDPOINT REDUCTION WITH FLUCTUATIONS
    \item \textbf{Asymptotic detailed balance:} NOT YET PROVED
    \item \textbf{Generation of Borel measure theory from the discrete model:} NOT CLAIMED
    \item \textbf{Physical / anyonic interpretation of divisor or winding data:} HEURISTIC / NOT CLAIMED
\end{itemize}



% -----------------------------------------------------------------------------
% Intended placement in the extended master architecture: Section 11.6.1.
% The compact V4 source does not yet contain the complete Chapter 11 hierarchy.
% Therefore the semantic number is retained in the heading without falsifying
% the automatic counters of this standalone source.
\subsection*{11.6.1 Quaternionic Three-Dimensional Contraction of the Fano Rotational Term}
\addcontentsline{toc}{subsection}{11.6.1 Quaternionic Three-Dimensional Contraction of the Fano Rotational Term}
\label{sec:fano-3d-contraction}

Let $\phi_{ijk}$ denote the totally antisymmetric Fano three-form associated with
an oriented octonionic multiplication table.  On any chosen quaternionic
subalgebra $\operatorname{span}\{1,e_a,e_b,e_c\}\cong\mathbb H$, the restriction
of $\phi$ to the imaginary three-plane is, after choosing a standard
Levi--Civita tensor $\epsilon$, of the form
\begin{equation}
\phi_{ijk}\big|_{\{a,b,c\}}=\sigma\,\epsilon_{ijk},\qquad \sigma\in\{+1,-1\}.
\end{equation}
The sign $\sigma$ records the relative orientation of the chosen Fano line and
the standard orientation of the physical three-space.  In particular, for the
orientation convention used in the present audit, the Fano line
$(1,3,2)$ gives
\begin{equation}
\phi_{123}=-1=-\epsilon_{123},
\end{equation}
so that $\sigma=-1$ on the selected $\{1,2,3\}$ sector.

Let the antisymmetric part of the three-dimensional velocity gradient be
\begin{equation}
\Omega_{jk}=\frac12(\partial_j u_k-\partial_k u_j)
          =-\frac12\epsilon_{jkl}\omega_l,
\end{equation}
with the displayed sign fixing the vorticity convention.  Define the restricted
Fano rotational contraction by
\begin{equation}
\mathcal R_F^{(3)}(\omega)
:=\sum_{i,j,k=1}^{3}\phi_{ijk}\,\omega_i\Omega_{jk}.
\end{equation}

\begin{theorem}[Orientation-dependent quaternionic contraction]
Under the above conventions,
\begin{equation}
\boxed{
\mathcal R_F^{(3)}(\omega)=-\sigma\,|\omega|^2.
}
\end{equation}
Hence the contraction is negative definite only when the Fano orientation and
the spatial orientation are aligned ($\sigma=+1$).  For the project convention
$(1,3,2)$, one has $\sigma=-1$ and therefore
\begin{equation}
\mathcal R_F^{(3)}(\omega)=+|\omega|^2.
\end{equation}
\end{theorem}

\begin{proof}
Using $\phi_{ijk}=\sigma\epsilon_{ijk}$ and the three-dimensional vorticity
identity,
\begin{align}
\mathcal R_F^{(3)}
&=-\frac{\sigma}{2}
  \sum_{i,j,k,l=1}^{3}
  \epsilon_{ijk}\epsilon_{jkl}\,\omega_i\omega_l \\
&=-\frac{\sigma}{2}
  \sum_{i,l=1}^{3}2\delta_{il}\,\omega_i\omega_l \\
&=-\sigma\sum_{i=1}^{3}\omega_i^2.
\end{align}
This proves the identity exactly.
\end{proof}

The preceding theorem is an algebraic contraction identity.  It does not, by
itself, provide a coercive estimate for the classical Navier--Stokes vortex
stretching term.  Two separate issues remain.

First, changing the relative orientation $\sigma$ changes the sign of the
contraction.  Consequently, negative definiteness is not an orientation-free
consequence of the Fano algebra.  A dissipative interpretation therefore
requires an independent dynamical or geometric principle that fixes the
physical embedding, the orientation, and the coupling sign.

Second, the classical three-dimensional enstrophy production is
\begin{equation}
\mathcal S_{\mathrm{NS}}
=\int_{\mathbb R^3}\omega_i S_{ij}\omega_j\,dx,
\end{equation}
where $S$ is the symmetric strain tensor.  The Fano contraction above is a
separate algebraic observable unless one proves that it enters the actual
Navier--Stokes evolution with a fixed coefficient and controls
$\mathcal S_{\mathrm{NS}}$ in a scaling-critical norm.  No such coercive bridge
is established by the contraction identity alone.

\paragraph{Gate status.}
The precise status of the sixth Fano gate is therefore
\[
\boxed{
\begin{array}{ll}
G_{\mathrm{NS}}^{(6a)} & \mathrm{EXACT\ PASS}\quad
  (\phi_{ijk}S_{jk}=0),\\
G_{\mathrm{NS}}^{(6b)}\text{ algebraic reduction} & \mathrm{EXACT\ PASS},\\
G_{\mathrm{NS}}^{(6b)}\text{ 3D contraction} &
  \mathcal R_F^{(3)}=-\sigma|\omega|^2,\\
\text{orientation-independent damping} & \mathrm{NOT\ ESTABLISHED},\\
G_{\mathrm{NS}}^{(6c)}\text{ coercive bridge} & \mathrm{OPEN},\\
\text{global 3D regularity} & \mathrm{NOT\ ESTABLISHED}.
\end{array}}
\]



% -----------------------------------------------------------------------------
% Intended placement in the extended master architecture: Section 11.6.4.
% The compact V4 source does not yet contain the complete Chapter 11 hierarchy.
% Therefore the semantic number is retained in the heading without falsifying
% the automatic counters of this standalone source.
\subsection*{11.6.4 Morrey Threshold and the Discrete $W^{1,6}$ Direction-Coherence Gate}
\addcontentsline{toc}{subsection}{11.6.4 Morrey Threshold and the Discrete W1,6 Direction-Coherence Gate}
\label{sec:fano-morrey-w16}

The Fano-angle identity on a selected quaternionic three-plane gives an exact
orientation-sensitive representation of the angle between two nonzero
vorticity directions.  Writing
\[
\xi(x,t)=\frac{\omega(x,t)}{|\omega(x,t)|}
\]
wherever $\omega\neq0$, and using
$\phi|_{\mathbb H_{\rm Im}}=\sigma\epsilon$ with $\sigma\in\{\pm1\}$, one has
\begin{equation}
\boxed{
\|\phi(\xi(x,t),\xi(y,t),\cdot)\|
=\|\xi(x,t)\times\xi(y,t)\|
=|\sin\theta(x,y,t)|.
}
\end{equation}
This identity is exact, but by itself it supplies no spatial modulus of
continuity.  The remaining problem is therefore analytic: a discrete or
coarse-grained EABC mechanism must yield a metric estimate strong enough to
control the spatial variation of $\xi$.

\paragraph{Sobolev obstruction in three dimensions.}
A quadratic lattice energy of Dirichlet type can at best suggest an
$H^1=W^{1,2}$ bound in the continuum.  In three dimensions the standard
Sobolev embedding gives
\begin{equation}
W^{1,2}(\mathbb R^3)\hookrightarrow L^6(\mathbb R^3),
\qquad
W^{1,2}(\mathbb R^3)\not\hookrightarrow C^{0,\alpha}(\mathbb R^3)
\quad(\alpha>0).
\end{equation}
Consequently, a purely quadratic direction-energy bound does not imply the
Hölder control required by geometric depletion criteria for vortex stretching.

\paragraph{Morrey threshold.}
For $p>3$, Morrey's inequality yields, locally on regular bounded domains (and
in the corresponding whole-space form),
\begin{equation}
W^{1,p}\hookrightarrow C^{0,\alpha},
\qquad
\alpha=1-\frac{3}{p}.
\end{equation}
Hence $p=6$ is the natural sufficient exponent for obtaining the one-half
Hölder modulus
\begin{equation}
\boxed{
W^{1,6}\hookrightarrow C^{0,1/2}.
}
\end{equation}
Equivalently, for a representative of $\xi$ on a suitable domain $\Omega$,
\begin{equation}
|\xi(x)-\xi(y)|
\le C_{\Omega}\|\nabla\xi\|_{L^6(\Omega)}|x-y|^{1/2}.
\end{equation}
This is a sufficient analytic route to a one-half Hölder estimate; it is not
claimed to be a necessary or unique route to Navier--Stokes regularity.

\paragraph{Scale-consistent discrete candidate.}
Let $h>0$ denote the lattice spacing and let $x\sim y$ denote nearest-neighbour
edges of a three-dimensional lattice.  The scale-consistent sixth-order
finite-difference energy is
\begin{equation}
\boxed{
E_h^{(6)}[\xi_h]
:=h^{-3}\sum_{x\sim y}|\xi_x-\xi_y|^6.
}
\end{equation}
Indeed, $|\xi_x-\xi_y|\sim h|\nabla\xi|$ and the number of edges per unit
volume is of order $h^{-3}$, so that the prefactor $h^{-3}=h^{3-6}$ is the
one compatible with a continuum $L^6$ gradient energy.  A uniform estimate
\begin{equation}
\sup_{0<h<h_0}E_h^{(6)}[\xi_h]\le C
\label{eq:w16-discrete-gate}
\end{equation}
would therefore be the appropriate discrete compactness input for a
$W^{1,6}$-type continuum bound, provided that the interpolation scheme,
domain control, and boundary conditions are specified and proved compatible.

\paragraph{A necessary sign-lift refinement of the Fano energy.}
The Fano norm measures $|\sin\theta|$, not the full vector difference.  For
unit vectors $a,b\in S^2$,
\begin{equation}
|a-b|^2=2(1-a\cdot b),
\qquad
|a\times b|^2=(1-a\cdot b)(1+a\cdot b).
\end{equation}
Thus $|a\times b|$ vanishes for both parallel and antiparallel vectors.  It
cannot, without an additional orientation condition, control $|a-b|$.  This
prevents one from inferring \eqref{eq:w16-discrete-gate} directly from a Fano
energy of the form
\begin{equation}
E_{h,F}^{(6)}
:=h^{-3}\sum_{x\sim y}
\|\phi(\xi_x,\xi_y,\cdot)\|^6
=h^{-3}\sum_{x\sim y}|\sin\theta_{xy}|^6.
\end{equation}
A local sign or hemisphere lock is sufficient.  For example, if neighbouring
directions satisfy
\begin{equation}
\xi_x\cdot\xi_y\ge0,
\label{eq:hemisphere-lock}
\end{equation}
then
\begin{equation}
|\xi_x-\xi_y|^2
=2(1-\xi_x\cdot\xi_y)
\le2|\xi_x\times\xi_y|^2,
\end{equation}
and hence
\begin{equation}
E_h^{(6)}[\xi_h]\le 8E_{h,F}^{(6)}[\xi_h].
\end{equation}
Therefore a genuine arithmetic-to-spatial bridge requires two independent
ingredients: a uniform sixth-order Fano-angle bound and a mechanism excluding
uncontrolled nearest-neighbour sign flips.

\begin{theorem}[Conditional discrete-to-Morrey bridge]
Let $\Omega\subset\mathbb R^3$ be a regular bounded domain and let $\xi_h$ be
unit-vector lattice fields admitting a stable interpolation
$I_h\xi_h\in W^{1,6}(\Omega)$.  Assume that the interpolation satisfies
\begin{equation}
\|\nabla I_h\xi_h\|_{L^6(\Omega)}^6
\le C_0 E_h^{(6)}[\xi_h]
\end{equation}
with $C_0$ independent of $h$, and that
\begin{equation}
\sup_h E_h^{(6)}[\xi_h]<\infty.
\end{equation}
Then the interpolants are uniformly bounded in $W^{1,6}(\Omega)$ up to the
usual control of their $L^6$ norm.  Morrey's inequality consequently gives
\begin{equation}
[I_h\xi_h]_{C^{0,1/2}(\Omega)}\le C
\end{equation}
with $C$ independent of $h$.  After extraction of a subsequence one obtains
weak $W^{1,6}$ compactness and, for every $\beta<1/2$, strong compactness in
$C^{0,\beta}$ on compact subsets under the standard hypotheses.
\end{theorem}

The theorem is deliberately conditional: the EABC framework has not yet
proved the uniform bound \eqref{eq:w16-discrete-gate}, the interpolation
estimate, or the hemisphere-lock condition \eqref{eq:hemisphere-lock}.  In
addition, $\xi=\omega/|\omega|$ is undefined on the zero-vorticity set; a
complete Navier--Stokes application must either work on high-vorticity regions
or establish a compatible extension procedure.

\paragraph{Relation to geometric depletion.}
The Constantin--Fefferman direction-of-vorticity program shows that sufficient
geometric coherence of vorticity directions can deplete the vortex-stretching
nonlinearity.  The role of the present subsection is narrower: it identifies a
concrete sufficient analytic target by which an EABC lattice mechanism could
produce a one-half Hölder modulus.  No equivalence with the full hypotheses of
any particular classical regularity theorem is asserted here.

\paragraph{Gate status.}
The resulting status is
\[
\boxed{
\begin{array}{ll}
G_{\mathrm{NS}}^{(8a)}\text{ topological-to-metric implication}
  & \mathrm{FAILED\ WITHOUT\ AN\ ENERGY\ ESTIMATE},\\
G_{\mathrm{NS}}^{(8b)}\text{ Sobolev obstruction}
  & W^{1,2}(\mathbb R^3)\not\hookrightarrow C^{0,\alpha},\\
G_{\mathrm{NS}}^{(8c)}\text{ Morrey threshold}
  & W^{1,6}\hookrightarrow C^{0,1/2}\quad\mathrm{STANDARD\ THEOREM},\\
G_{\mathrm{NS}}^{(8d)}\text{ uniform discrete }W^{1,6}\text{ bound}
  & \mathrm{OPEN},\\
G_{\mathrm{NS}}^{(8e)}\text{ Fano-to-vector sign lift}
  & \mathrm{OPEN\ WITHOUT\ ORIENTATION\ LOCK},\\
G_{\mathrm{NS}}^{(8f)}\text{ arithmetic-to-spatial bridge}
  & \mathrm{OPEN},\\
\text{global 3D Navier--Stokes regularity}
  & \mathrm{NOT\ ESTABLISHED}.
\end{array}}
\]


% -----------------------------------------------------------------------------
% Intended placement in the extended master architecture: Section 11.6.5.
% The compact V4 source does not yet contain the complete Chapter 11 hierarchy.
% Therefore the semantic number is retained in the heading without falsifying
% the automatic counters of this standalone source.
\subsection*{11.6.5 Quartic $E_8$ Root Isotropy and a Conditional $H^2$ UV-Coercivity Route}
\addcontentsline{toc}{subsection}{11.6.5 Quartic E8 Root Isotropy and a Conditional H2 UV-Coercivity Route}
\label{sec:e8-quartic-uv}

The direct $W^{1,6}$ route of Section~\ref{sec:fano-morrey-w16} asks for a
sixth-order discrete gradient bound.  A spectrally more natural sufficient
route is to obtain quadratic control of two spatial derivatives.  In three
space dimensions, local $H^2$ control (together with the standard lower-order
localization terms) implies
\begin{equation}
H^2(\Omega)\hookrightarrow W^{1,6}(\Omega)
\hookrightarrow C^{0,1/2}(\Omega)
\end{equation}
on regular bounded domains.  Thus an isotropic positive fourth-order Fourier
symbol is a natural candidate for a direction-coherence mechanism.

\paragraph{Exact fourth moment of the $E_8$ root system.}
Let $R(E_8)\subset\mathbb R^8$ denote the standard set of $240$ roots,
normalized by $|r|^2=2$.  The normalized root directions form a spherical
$7$-design.  Consequently all spherical moments through degree four coincide
with their rotationally invariant spherical averages.  In particular, for
arbitrary $k\in\mathbb R^8$,
\begin{theorem}[Quartic $E_8$ root-moment identity]
\begin{equation}
\boxed{
\sum_{r\in R(E_8)}(r\cdot k)^4=36\,|k|^4.
}
\label{eq:e8-fourth-moment}
\end{equation}
\end{theorem}

\begin{proof}
Write $u=r/\sqrt2\in S^7$.  Since the normalized roots are a spherical
$7$-design, their fourth moment equals the fourth spherical moment.  For the
unit sphere in dimension $d=8$,
\begin{equation}
\frac{1}{|S^7|}\int_{S^7}(u\cdot k)^4\,dS(u)
=\frac{3}{d(d+2)}|k|^4
=\frac{3}{80}|k|^4.
\end{equation}
Because $(r\cdot k)^4=4(u\cdot k)^4$ and there are $240$ roots,
\begin{equation}
\sum_{r\in R(E_8)}(r\cdot k)^4
=240\cdot4\cdot\frac{3}{80}|k|^4
=36|k|^4.
\end{equation}
\end{proof}

Equivalently, the fourth moment tensor is
\begin{equation}
\sum_{r\in R(E_8)}r_i r_j r_k r_l
=12\bigl(\delta_{ij}\delta_{kl}
        +\delta_{ik}\delta_{jl}
        +\delta_{il}\delta_{jk}\bigr).
\label{eq:e8-fourth-tensor}
\end{equation}
Equation~\eqref{eq:e8-fourth-moment} is therefore an exact algebraic-geometric
identity.  It is logically independent of the dimensional pole of the
$E_8$ Epstein zeta function at $s=8$; the latter supplies spectral information
about the lattice but does not by itself generate the fourth-order operator
below.

\paragraph{Dimensionally consistent three-dimensional pullback.}
A Navier--Stokes application lives in physical dimension three, whereas the
roots $r$ lie in $\mathbb R^8$.  Therefore the expression $r\cdot\nabla$ is not
defined until an embedding is specified.  Let
\begin{equation}
\iota:\mathbb R^3\longrightarrow\mathbb R^8,
\qquad \iota^T\iota=I_3,
\end{equation}
be an isometric embedding of the physical spatial cotangent space, and define
the pulled-back root directions
\begin{equation}
q_r:=\iota^T r\in\mathbb R^3.
\end{equation}
For every physical wave vector $\kappa\in\mathbb R^3$, applying
\eqref{eq:e8-fourth-moment} to $k=\iota\kappa$ gives the exact identity
\begin{equation}
\boxed{
\sum_{r\in R(E_8)}(q_r\cdot\kappa)^4
=36|\kappa|^4.
}
\label{eq:e8-pullback-fourth}
\end{equation}
Thus the quartic isotropy survives every isometric choice of physical
three-plane.  No preferred three-plane is selected by this identity alone.

\paragraph{Exact quadratic fourth-order form.}
For a sufficiently regular vector field $\xi:\mathbb R^3\to\mathbb R^3$ with
suitable decay, define
\begin{equation}
\mathcal Q_{E_8}[\xi]
:=\sum_{r\in R(E_8)}
\bigl\|(q_r\cdot\nabla)^2\xi\bigr\|_{L^2}^2.
\end{equation}
By Plancherel and \eqref{eq:e8-pullback-fourth}, componentwise in $\xi$,
\begin{equation}
\boxed{
\mathcal Q_{E_8}[\xi]
=36\int_{\mathbb R^3}|\kappa|^4|\widehat\xi(\kappa)|^2\,d\kappa
=36\,\|D^2\xi\|_{L^2}^2,
}
\label{eq:e8-h2-form}
\end{equation}
where the last equality uses the standard Frobenius norm of the Hessian.  This
is an exact isotropic fourth-order quadratic form.

\begin{theorem}[Conditional $E_8$ UV-coercivity bridge]
Assume that the effective EABC evolution or action for the vorticity direction
contains, on a specified physical three-plane, a positive term
\begin{equation}
\gamma\,\mathcal Q_{E_8}[\xi],
\qquad \gamma>0,
\label{eq:e8-positive-gamma}
\end{equation}
and that the resulting a priori energy estimate controls this term uniformly
in time together with the lower-order norms required for localization.  Then
\eqref{eq:e8-h2-form} yields uniform local $H^2$ control of $\xi$.  On regular
bounded domains this implies
\begin{equation}
\xi\in H^2
\quad\Longrightarrow\quad
\xi\in W^{1,6}
\quad\Longrightarrow\quad
\xi\in C^{0,1/2},
\end{equation}
with the corresponding norm bounds.
\end{theorem}

The theorem is conditional in an essential sense.  The exact root identity
fixes the isotropic symbol once the quadratic form $\mathcal Q_{E_8}$ is
present, but it does not prove that the Navier--Stokes evolution contains such
a term, does not determine a positive coupling $\gamma$, and does not establish
an a priori bound for that term.  Moreover, since $\xi=\omega/|\omega|$ is
undefined on the zero-vorticity set and need not lie in global $L^2(\mathbb
R^3)$, a complete application must be formulated locally on high-vorticity
regions or with a justified extension/localization procedure.

\paragraph{Relation to the spectral pole at $s=8$.}
The dimensional pole of the radial $E_8$ Epstein zeta function at $s=8$ and
the fourth-root moment identity encode different pieces of structure.  The
pole is a meromorphic spectral invariant tied to dimension and covolume.  The
identity \eqref{eq:e8-fourth-moment} is a finite root-design identity.  A
spectral-action derivation connecting the two and producing the positive
coefficient $\gamma$ in \eqref{eq:e8-positive-gamma} is not established here.
This separation prevents the pole itself from being interpreted as an
unproved source of ultraviolet coercivity.

\paragraph{Gate status.}
The ninth Navier--Stokes symmetry gate is therefore
\begin{center}
\begin{boxedminipage}{0.96\textwidth}
\small
\begin{tabularx}{\linewidth}{@{}lX@{}}
$G_{\mathrm{NS}}^{(9a)}$ dimensional $E_8$ pole at $s=8$
  & \textbf{EXACT/STANDARD}.\\
$G_{\mathrm{NS}}^{(9b)}$ quartic root isotropy
  & $\sum_r(r\cdot k)^4=36|k|^4$; \textbf{EXACT PASS}.\\
$G_{\mathrm{NS}}^{(9c)}$ 3D isometric pullback
  & $\sum_r(q_r\cdot\kappa)^4=36|\kappa|^4$; \textbf{EXACT PASS}.\\
$G_{\mathrm{NS}}^{(9d)}$ positive spectral-action coupling $\gamma$
  & \textbf{OPEN}.\\
$G_{\mathrm{NS}}^{(9e)}$ $H^2$ coercivity
  & \textbf{CONDITIONAL} on $\gamma>0$ and a uniform a priori energy bound.\\
$G_{\mathrm{NS}}^{(9f)}$ $H^2\to W^{1,6}\to C^{0,1/2}$
  & \textbf{STANDARD ANALYSIS}.\\
$G_{\mathrm{NS}}^{(9g)}$ vorticity-direction regularity bridge
  & \textbf{CONDITIONAL}.\\
Global 3D Navier--Stokes regularity
  & \textbf{NOT ESTABLISHED}.\\
\end{tabularx}
\end{boxedminipage}
\end{center}


\end{document}
